\section{Pricing on the lattice}
\label{sec:lvlattice}

\subsection{The implicit march and its M-matrix}

The forward equation must now be solved numerically, and this section
chooses the scheme with one criterion above all others: it must
inherit \cref{prop:lvpropagate}, not merely approximate it.

First the setting.  The equation is solved on a strike grid over
$[0,y_{\max}]$, with the two end values pinned: $c(\cdot,0)=1$ (a call
struck at zero delivers the forward, worth $1$ in normalized units)
and $c(\cdot,y_{\max})=0$ (far enough out, the call is worthless).
Pinned end values are called \emph{Dirichlet} boundaries.  The march
then steps maturity forward in increments $\Delta\tau$, and at each
step it must decide at which time level to evaluate the strike
derivative.  The \emph{fully implicit} choice --- evaluate at the new
level --- is the one this chapter defends.

On the grid, the second derivative becomes a weighted second
difference, and the weights are worth deriving because their signs
carry the whole theory.  The grid need not be uniform; write
$\Delta y_i^-=y_i-y_{i-1}$ and $\Delta y_i^+=y_{i+1}-y_i$ for the gaps
around node $i$.  Taylor-expand the two neighbors of a smooth function
$c$ around $y_i$:
\[
c_{i\pm1}=c_i\pm\Delta y_i^{\pm}\,c'(y_i)
+\tfrac12\,(\Delta y_i^{\pm})^{2}\,c''(y_i)+O(\Delta y^{3}).
\]
Multiply the ``$+$'' expansion by $\Delta y_i^-$, the ``$-$''
expansion by $\Delta y_i^+$, and add: the first-derivative terms
cancel, leaving
\[
\Delta y_i^-\,c_{i+1}+\Delta y_i^+\,c_{i-1}
-(\Delta y_i^-+\Delta y_i^+)\,c_i
=\tfrac12\,\Delta y_i^+\Delta y_i^-(\Delta y_i^++\Delta y_i^-)\,
c''(y_i)+O(\Delta y^{4}),
\]
and solving for $c''$ gives the three-point weights.  Multiplying by
$\tfrac12 v_i y_i^{2}$, the discrete generator $\mathcal L$ encoding
$\tfrac12 v y^{2}\partial_{yy}$ has the stencil
\begin{equation}
\mathcal L_{i,i\mp1}
=\frac{v_i\,y_i^{2}}{(\Delta y_i^-+\Delta y_i^+)\,\Delta y_i^{\mp}}\;\ge0,
\qquad
\mathcal L_{ii}=-\big(\mathcal L_{i,i-1}+\mathcal L_{i,i+1}\big),
\label{eq:lvstencil}
\end{equation}
which reduces to the familiar $(1,-2,1)/\Delta y^{2}$ pattern on a
uniform grid.  Note the two signs that will do all the work: the
off-diagonal entries are nonnegative, and each diagonal is exactly
minus the sum of its row's off-diagonals.

An implicit step then solves a linear system for the vector $U^{n+1}$
of calls at the new maturity, the operator evaluated at the new level:
\begin{equation}
\mathcal M^{n+1}U^{n+1}=U^{n}+(\text{boundary terms}),
\qquad
\mathcal M^{n+1}=I-\Delta\tau\,\mathcal L^{n+1}.
\label{eq:lvstep}
\end{equation}
Because the stencil couples only neighbors, $\mathcal M$ is
tridiagonal, and each step is one tridiagonal solve --- linear cost in
the number of nodes.  The boundary terms carry the pinned values $1$
and $0$ through the stencil's end columns, and two later proofs lean
on them, so write them out once as a vector $b^{n+1}$: only the row
whose lower neighbor is the boundary node $y=0$ has a nonzero entry,
\[
b^{n+1}_i=\Delta\tau\,\mathcal L^{n+1}_{i,i-1}\times1
\quad(\text{the pinned call value at }y=0),
\qquad
b^{n+1}_i=0\ \text{ at every other row}
\]
--- the far boundary would contribute through its own stencil arm, but
its pinned value is $0$.

Now the promised discipline.  \Cref{prop:lvpropagate} was proved for
the continuous equation.  Does the discrete march \emph{inherit} it
--- no invented negative densities, no broken calendar order --- or
merely approximate it, with errors of unknown sign?  The answer runs
through one classical notion from linear algebra.  A matrix with
positive diagonal, nonpositive off-diagonal entries and enough
diagonal dominance is called an \emph{M-matrix}, and the property that
sign pattern buys is exactly the one needed here: the inverse of an
M-matrix has no negative entries.

\begin{proposition}[The step matrix is an M-matrix]
\label{prop:lvmmatrix}
For any $v\ge0$ and any $\Delta\tau>0$, the matrix
$\mathcal M=I-\Delta\tau\mathcal L$ of \eqref{eq:lvstep} has positive
diagonal, non-positive off-diagonals, and each diagonal entry exceeds the
absolute sum of its row's off-diagonals by at least $1$ (exactly $1$ at
fully interior rows).  Hence $\mathcal M$ is a strictly diagonally
dominant nonsingular M-matrix, and $\mathcal M^{-1}\ge0$ elementwise.
\end{proposition}

\begin{proof}
The sign pattern reads off \eqref{eq:lvstencil}: the off-diagonals of
$\mathcal M$ are $-\Delta\tau\mathcal L_{i,i\mp1}\le0$; the diagonal is
$1+\Delta\tau(\mathcal L_{i,i-1}+\mathcal L_{i,i+1})\ge1$; and the
dominance gap is the $1$ contributed by the identity at interior rows,
more at rows whose stencil arm reaches a boundary column (that
neighbor is not among the unknowns).  This is the defining pattern of
a nonsingular M-matrix, and such matrices are inverse-positive
\citep{BermanPlemmons1994}.  A self-contained argument fits in four
lines.  To solve $\mathcal M u=b$ with $b\ge0$, split
$\mathcal M=D-\mathcal N$ into its diagonal $D>0$ and off-diagonal
part $\mathcal N\ge0$, and iterate Jacobi's method from zero:
\[
u^{(m+1)}=D^{-1}\big(\mathcal N u^{(m)}+b\big),\qquad u^{(0)}=0 .
\]
Every iterate is nonnegative, being assembled from nonnegative pieces;
the iteration converges, because every row of $D^{-1}\mathcal N$ sums
to less than one (that is diagonal dominance), so each update
multiplies the largest error entry by a factor below one; and the
limit solves $\mathcal M u=b$.  So
$b\ge0$ forces $u\ge0$, for every such $b$ --- which is the statement
$\mathcal M^{-1}\ge0$, column by column.
\end{proof}

Inverse positivity is the discrete counterpart of heat flow never
turning a nonnegative temperature profile negative, and it translates
directly into the property numerical
analysts call \emph{monotonicity} --- order in, order out --- and into
stability in the strongest everyday sense: the largest absolute value
in the solution can never grow.

\begin{corollary}[Discrete maximum principle, unconditionally]
\label{cor:lvdmp}
For any $\Delta\tau>0$ the scheme \eqref{eq:lvstep} is \emph{monotone} ---
$U^{n}\ge V^{n}$ with the same boundary data implies $U^{n+1}\ge V^{n+1}$
--- and every step keeps the values between the extremes of the previous
step and the boundary data: the march is unconditionally
$\ell^{\infty}$-stable (the maximum absolute value never grows), with
no step-size restriction, and cannot create
new extrema.
\end{corollary}

\begin{proof}
Monotonicity is inverse positivity applied to a difference: if
$U^{n}\ge V^{n}$, the two systems \eqref{eq:lvstep} subtract to
$\mathcal M^{n+1}(U^{n+1}-V^{n+1})=U^{n}-V^{n}\ge0$, and multiplying
by $\mathcal M^{-1}\ge0$ preserves the sign.  For the upper bound, let
$U_\ast$ be the larger of $\max_iU^{n}_i$ and the boundary values $1$
and $0$, and test the constant vector $U_\ast\1$, where $\1$ is the
all-ones vector (a vector here, not the indicator of
\cref{prop:lvforward}'s proof).  At a fully interior row, the row of
$\mathcal M$ sums to $1$ (because the stencil's row sums to zero), so
$(\mathcal M\,U_\ast\1)_i=U_\ast\ge U^{n}_i$.  At a row next to a
boundary, the arm that would reach the boundary column is missing from
$\mathcal M$, leaving the larger value
$(\mathcal M\,U_\ast\1)_i=U_\ast+\Delta\tau\,\mathcal L_{i,\rm bd}\,U_\ast$,
with $\mathcal L_{i,\rm bd}$ the missing arm's weight; the right-hand
side there is $U^{n}_i+b^{n+1}_i$ with
$b^{n+1}_i\le\Delta\tau\,\mathcal L_{i,\rm bd}\times1$.  Since
$U_\ast$ dominates both $U^{n}_i$ and the boundary value $1$, the left
side dominates the right at every row:
$\mathcal M(U_\ast\1)\ge U^{n}+b^{n+1}$ entrywise.  Multiplying by
$\mathcal M^{-1}\ge0$ gives $U_\ast\1\ge U^{n+1}$.  The lower bound is
symmetric.
\end{proof}

\begin{remark}[What the M-matrix argument does and does not prove]
\label{rem:lvscope}
\Cref{prop:lvmmatrix,cor:lvdmp} control the call \emph{values}: order
preservation, boundedness, no invented extrema.  They do not by
themselves give nonnegative second differences --- a bounded monotone
vector can still be locally concave --- so ``the scheme cannot
manufacture negative densities'' would need a separate
convexity-preservation proof for the actual stencil and boundary
treatment, and this chapter does not claim one.  What it does instead
is measure: the scheme is monotone and stable by theorem, and the
butterfly and calendar cleanliness of the \emph{output} is checked on
the price grid with every fit.  \Cref{sec:lvexamples} reports the
measured minima for the chapter's two calibrated sheets; they sit at
the scale of solver rounding.
\end{remark}

\subsection{Why not Crank--Nicolson}

The implicit march is only first-order accurate in the time step, and
the standard second-order upgrade is Crank--Nicolson: evaluate the
operator half at the old level and half at the new, replacing
\eqref{eq:lvstep} by
\[
\big(I-\tfrac{\Delta\tau}{2}\mathcal L\big)U^{n+1}
=\big(I+\tfrac{\Delta\tau}{2}\mathcal L\big)U^{n}.
\]
Look at what the upgrade costs.  The implicit factor on the left is
still an M-matrix --- that part survives.  But monotonicity now also
needs the \emph{explicit} factor on the right to map nonnegative
vectors to nonnegative vectors, and its diagonal entries are
$1-\tfrac{\Delta\tau}{2}(\mathcal L_{i,i-1}+\mathcal L_{i,i+1})$:
nonnegative only when the step is small enough,
\begin{equation}
\Delta\tau\;\le\;\frac{2}{\mathcal L_{i,i-1}+\mathcal L_{i,i+1}}
\;=\;\frac{2\,\Delta y^{2}}{v_i\,y_i^{2}}
\quad\text{(uniform grid)}.
\label{eq:lvcfl}
\end{equation}
This is a step-size restriction of the Courant--Friedrichs--Lewy (CFL)
type --- the generic fate of schemes with an explicit part --- and it
tightens with the \emph{square} of
the strike step.  Put numbers in: at a $40\%$ local volatility on a
$\Delta y=0.005$ lattice --- the resolution \cref{sec:lvinfo} will
show short-dated smiles require --- the bound is
$\MacLvLatCflBound$ years, more than thirty times smaller than the
standard $0.01$ marching step.  Monotone Crank--Nicolson on realistic
lattices is simply not on offer.  When the bound is violated, the
scheme remains second-order \emph{accurate} but is no longer
\emph{monotone}, and the payoff kink at $y=1$ seeds the oscillation
that \cref{fig:lvmonotone} displays.

\begin{figure}[htbp]
\centering
\paperfig{0.9\textwidth}{fig_lv_monotone.pdf}
\caption{The same flat $40\%$ surface, the same $\Delta y=0.005$
lattice, marched to $\tau=0.1$ in ten steps by the same solver ---
only the scheme differs.  (a)~Implicit Euler: the discrete density is
a clean bell, minimum \MacLvLatImplMinDens\ in the plotted window.
(b)~Undamped Crank--Nicolson: the kink oscillation swings the discrete
density to \MacLvLatCnMinDens.}
\label{fig:lvmonotone}
\end{figure}

\Cref{fig:lvmonotone} shows what breaking monotonicity does to the
object we actually care about: the discrete density, i.e.\ the second
differences of the marched prices.  Both panels march the same flat
$40\%$ surface on the same $\Delta y=0.005$ lattice to $\tau=0.1$ in
ten steps; only the scheme differs.  In panel~(a), implicit Euler: the
second differences form a clean bell --- a discrete Gaussian, as they
should --- whose minimum over the plotted window is
\MacLvLatImplMinDens: nonnegative, exactly as \cref{cor:lvdmp}
guarantees.  In panel~(b), undamped Crank--Nicolson: the oscillation
seeded by the payoff kink swings the discrete density down to
\MacLvLatCnMinDens\ --- a negative density, which is butterfly
arbitrage, manufactured by the scheme that is, in the textbook
ordering, the more accurate of the two.  Damped start-up steps --- a
few implicit steps before switching to Crank--Nicolson --- are the
standard repair \citep{GilesCarter2006}.  The reference implementation
declines the trade altogether: it pays one order of time accuracy and
keeps \cref{cor:lvdmp} unconditionally.

\subsection{Reading out, and the operator's blind spot}

The observation operator is plain: the normalized call at a quote's
$(\tau,y)$ is read off the marched solution by interpolation in
strike, and quoted expiries are forced to be lattice time-nodes, so no
interpolation in maturity is ever needed.

One principle from inverse problems completes the section, because it
dictates part of the chapter's protocol.  An optimizer facing a
\emph{fixed} discretization will happily bend its parameters to cancel
the operator's own discretization error.  The residual measured on the
fitting operator therefore flatters the fit, and validating a
reconstruction with the operator that produced it has a name in the
inverse-problems literature: the \emph{inverse crime}.  The chapter's
protocol answers it by repricing every fit on a \emph{refined}
operator --- half the strike step, four times the time steps --- and
reporting both numbers.  The gap between them is not noise; it is the
measured size of the compensation, and \cref{sec:lvexamples} will show
it is the dominant error term at the short end.
